\section{P\'olya's papers on zeros, one by one}\label{pol:sec:papers}

\subsection*{P1. Approximation by polynomials with real zeros (1913)}
\emph{\"Uber Ann\"aherung durch Polynome mit lauter reellen Wurzeln}, Rend.\ Circ.\ Mat.\ Palermo 36 \cite{Polya13}. Completing Laguerre, P\'olya characterised the entire functions that are locally uniform limits of real polynomials with only real zeros: they are exactly
\begin{equation}\label{pol:eq:LP}
\begin{gathered}
f(z)=cz^me^{-\lambda z^2+\kappa z}\prod_k\left(1-\frac{z}{b_k}\right)e^{z/b_k},\\[4pt] m\in\{0,1,2,\dots\},\ \lambda\ge0,\ c,\kappa\in\mathbb R,\ b_k\in\mathbb R\setminus\{0\},\ \sum_kb_k^{-2}<\infty,
\end{gathered}
\end{equation}
the Laguerre P\'olya class $\mathcal{LP}$. The functions of type I, $cz^me^{\sigma z}\prod_k(1+z/b_k)$ with $c\ge0$, $\sigma\ge0$, $b_k>0$ and $\sum_kb_k^{-1}<\infty$, are exactly the locally uniform limits of polynomials with only negative zeros and positive coefficients. RH is the statement $\Xi\in\mathcal{LP}$.

\subsection*{P2. P\'olya and Schur, two kinds of factor sequences (1914)}
\emph{\"Uber zwei Arten von Faktorenfolgen in der Theorie der algebraischen Gleichungen}, J.\ reine angew.\ Math.\ 144 \cite{PS14}. A sequence $(g_k)$ is a multiplier sequence if $\sum p_kx^k$ with only real zeros implies $\sum g_kp_kx^k$ with only real zeros. P\'olya and Schur proved that a positive sequence is a multiplier sequence if and only if $\sum g_kz^k/k!$ is in $\mathcal{LP}$ of type I, if and only if every Jensen polynomial $\sum_{k=0}^d\binom dkg_kx^k$ has only real zeros.

\subsection*{P3. Genus zero and one (1915)}
\emph{Algebraische Untersuchungen \"uber ganze Funktionen vom Geschlechte Null und Eins}, J.\ reine angew.\ Math.\ 145 \cite{Polya15}. Operations on entire functions of genus at most one that preserve reality of zeros. The statement used later is the shift lemma, quoted as Hilfssatz~II of the 1926 paper P8 in Section~\ref{rz:sec:polya}: if $G$ is real, of genus $0$ or $1$, with only real zeros, at least one of them, then $G(z+ia)+G(z-ia)$ has only real zeros for real $a$.

\subsection*{P4. Zeros of certain entire functions (1918)}
\emph{\"Uber die Nullstellen gewisser ganzer Funktionen}, Math.\ Z.\ 2 \cite{Polya18}. The continuous Enestr\"om Kakeya theorem: if $f$ is positive and nondecreasing on $(0,1)$, then $\int_0^1f(t)\cos(zt)\,dt$ and $\int_0^1f(t)\sin(zt)\,dt$ have only real zeros. The Riemann kernel is positive but decreasing on $(0,\infty)$, so the theorem does not apply.

\subsection*{P5. Number theoretic remarks (1919)}
\emph{Verschiedene Bemerkungen zur Zahlentheorie}, Jahresber.\ DMV 28 \cite{Polya19}. Among other things, the P\'olya conjecture that the summatory Liouville function $L(x)=\sum_{n\le x}\lambda(n)$ satisfies $L(x)\le0$ for $x\ge2$. It implies RH and is false: Haselgrove (1958) \cite{Haselgrove58} disproved it and Tanaka (1980) \cite{Tanaka80} found the first sign change at $x=906{,}150{,}257$.

\subsection*{P6. Zeros of successive derivatives (1922)}
\emph{\"Uber die Nullstellen sukzessiver Derivierten}, Math.\ Z.\ 12 \cite{Polya22}. P\'olya introduced the final set of a function, the set of limit points of the zeros of $f^{(n)}$ as $n\to\infty$, and determined it in several classes. The statements about real entire functions that we use are recorded under P13.

\subsection*{P7. A Fourier integral (1923)}
\emph{On the zeros of an integral function represented by Fourier's integral}, Messenger of Math.\ 52 \cite{Polya23m}. A short note on the zeros of cosine transforms of positive kernels, a precursor of the 1926 and 1927 papers.

\subsection*{P8. The integral representation of the Riemann xi function (1926)}
\emph{Bemerkung \"uber die Integraldarstellung der Riemannschen $\xi$ Funktion}, Acta Math.\ 48 \cite{Polya}. As $u\to+\infty$ the kernel is dominated by its $n=1$ term, and P\'olya replaced it by the even function
\begin{equation}\label{pol:eq:phistar}
\Phi^*(u)=4\pi^2\left(e^{9u/2}+e^{-9u/2}\right)e^{-\pi(e^{2u}+e^{-2u})}=8\pi^2\cosh\!\left(\tfrac{9u}{2}\right)e^{-2\pi\cosh 2u},
\end{equation}
and proved that $\Xi^*(t)=2\int_0^\infty\Phi^*(u)\cos(tu)\,du$ has only real zeros, with a zero count asymptotic to that of $\Xi$. The proof combines the reality of the zeros of $z\mapsto K_{iz}(a)$, $a>0$, with the shift lemma P3. A refinement keeping the second exponential term is also real rooted.

\subsection*{P9. Zeros of certain trigonometric integrals (1926)}
J.\ London Math.\ Soc.\ 1 \cite{Polya26jlms}. Following Hardy's proof of infinitely many zeros on the critical line, P\'olya explained why infinitely many real zeros is much easier than no complex zeros, and gave a general criterion for infinitely many real zeros. In Thorin language the infinitude of real zeros is the statement that the on line part $\varphi_{\mathrm{on}}=\prod_{\text{on line}}(1+s/\gamma^2)^{-1}$, a GGC whatever the zero set, has a Thorin measure with infinitely many atoms.

\subsection*{P10. Trigonometric integrals with only real zeros (1927)}
\emph{\"Uber trigonometrische Integrale mit nur reellen Nullstellen}, J.\ reine angew.\ Math.\ 158 \cite{Polya27}. The central paper. A real analytic $\phi$ is a universal factor if $\int\mathcal F(t)e^{izt}dt$ with only real zeros implies the same for $\int\phi(t)\mathcal F(t)e^{izt}dt$, for every admissible kernel $\mathcal F$, that is, every real kernel that decays fast enough for these integrals to define entire functions of $z$. P\'olya proved that $\phi$ is a universal factor if and only if $z\mapsto\phi(iz)$ belongs to $\mathcal{LP}$. Hence $e^{\lambda t^2}$ ($\lambda>0$) and $\cosh(bt)$ ($b\in\mathbb R$) are universal factors. He also proved a Mellin criterion: if $\mathcal F$ is admissible, holomorphic near the origin, and $H(z)=\int_0^\infty\mathcal F(t)t^{z-1}dt$ has no zeros in $\mathbb C\setminus(-\infty,0]$, then $\int\mathcal F(t^{2q})e^{izt}dt$ has only real zeros for every positive integer $q$. Consequences: $\int e^{-t^{2q}}e^{izt}dt$ and $\int\cosh(bt)e^{-a\cosh t}e^{izt}dt$ with $a>0$ and $b\in\mathbb R$ have only real zeros. The factor $e^{\lambda t^2}$ is the starting point of de Bruijn (1950) and Newman (1976), hence of the de Bruijn Newman constant $\Lambda$ (Table~\ref{tab:ch21:hist}).

\subsection*{P11. Jensen's investigations (1927)}
\emph{\"Uber die algebraisch funktionentheoretischen Untersuchungen von J.\ L.\ W.\ V.\ Jensen}, Kgl.\ Danske Vid.\ Selsk.\ Math.\ Fys.\ Medd.\ 7 \cite{Polya27j}. P\'olya edited and proved Jensen's criteria: a real entire function of genus at most one is in $\mathcal{LP}$ if and only if all its Jensen polynomials have only real zeros, together with Jensen's integral criteria. This is the source of the Jensen polynomial approach to RH (Section~\ref{fm:sec:polyanec}); Farmer (2022) \cite{Farmer22} argued that this direction cannot reach RH.

\subsection*{P12. Fourier's work on transcendental equations (1930)}
Quart.\ J.\ Math.\ Oxford 1 \cite{Polya30}. The Fourier P\'olya conjectures on the number of nonreal zeros of successive derivatives.

\subsection*{P13. Derivatives and analytic character (1943)}
Bull.\ AMS 49 \cite{Polya43}. The P\'olya Wiman conjecture: if $f$ is real entire of order below two with finitely many nonreal zeros, then $f^{(n)}$ has only real zeros for all large $n$; proved by Craven, Csordas and Smith (1987) \cite{CCS87} and Kim (1990) \cite{Kim90}, with the Fourier P\'olya conjecture settled by Ki and Kim (2000) \cite{KiKim00}. For $\Xi$, whose zeros lie in a strip, the theorem of Ki and Kim gives the local statement that for every $R$ the zeros of $\Xi^{(n)}$ in $|z|<R$ are real for all large $n$; Conrey showed that the proportion of real zeros of $\Xi^{(n)}$ tends to one \cite{Conrey83}. Global real rootedness of $\Xi^{(n)}$ for some $n$ is not known.

\subsection*{P14. The Hilbert P\'olya suggestion}
In a 1982 letter to Odlyzko, P\'olya recalled suggesting to Landau, around 1914, that the zeros might be the eigenvalues of a self adjoint operator. The Krein string route of Corollary~\ref{bp:cor:kent}, which assumes RH and simple zeros, is a probabilistic form of this.

